\scanpage{62}
sur une même $D$ il y a une structure
polygonale induite (cela a un sens \inserted{clair} si
le nb $s_D$ de ces sommets est $\geq3$). \inserted{L'ens.} des
orientations de $D$ (ou encore, l'ens. $\pi_0(\widetilde X\setminus\widetilde D)$)
est en corr. 1-1 avec l'ens. des orientations
\deleted{\unsure{de la}} de ce polygône combinatoire. \deleted{(Si $s_D=2$ ou $s_D=1$}
% slanted
\margin{Il vaut peut être mieux dire que $D$ est le support d'un 1-sous-complexe du squelette $K^1$ de dim 1 de la carte et ce sous-complexe est un polygône comb. d'ordre $s_D$ (sans exclure les cas $s_D=1$ ou $s_D=2$\ldots)}

Soit $s$ un sommet de la carte. Si on exclut le
cas où toutes les psdr. sont ``concourrantes'' en $s$
(auquel cas $s$ est l'unique sommet), il
existe donc une $D$ qui ne passe pas par $s$,
et l'ens \inserted{$\Omega_s(\unsure{x|X})$} des orientations de $X$ en $s$ est en
corr 1-1 can. avec l'ens \inserted{$\Omega(D)$} des orientations
du polygône comb. $D$ (qui est défini en
termes du 1-complexe $K^1$ et de sa décomposition en psdroites\ldots) Mais soit $D'$ \deleted{\unsure{et $D''$}}
une autre psdroite ne passant pas par $s$, comment
expliciter la bijection composée
\diagram{\begin{tikzcd}[ampersand replacement=\&,column sep=small,row sep=small]
\Omega(D)\arrow[rr,dash,"\sim"]\arrow[dr,"\sim"'] \& \& \Omega(D')\\
\& \Omega_s(x)\arrow[ur,"\sim"] \&
\end{tikzcd}}
? \deleted{Si par $s$} Comme $s$ est d'ordre $\geq4$,
il passe par $s$ deux psdroites $\Delta,\Delta'$, \deleted{\unsure{soient}} \inserted{supposons}
qu'on puisse \deleted{\unsure{les}} choisir $\Delta$ \uline{et} $\Delta'$ non \unsure{concourrantes|concourantes} avec $D,D'$, soient
$u,u'$ (resp. $v,v'$) les int. de $\Delta$ (resp. $\Delta'$) avec $D,D'$,
dont l'int est désignée par $t$. On constate alors
(puisque la configuration des quatre psdroites
% archive sheet number: 61
